\begin{abstract}
Safe physical human--robot interaction on floating-base robots is
fundamentally a problem of regulating interaction dynamics under changing
contact constraints.
This paper develops a contact-consistent normalization of those dynamics:
after priority-consistent cancellation of balance and contact tasks, the
end-effector interaction channel reduces to a linear double integrator whose
discrete state matrix is constant within each contact mode.
Robot configuration and support conditions enter only through the
contact-consistent task inertia and hence the input matrix.
This representation turns floating-base pHRI from a nonlinear whole-body
control problem into a configuration-invariant predictive interaction problem
with reusable rollouts, cached contact-mode matrices, and ${\geq}1$\,kHz
updates.
The resulting controller regulates the normalized interaction dynamics with a
receding-horizon QP, a contact-mode-indexed disturbance observer for
offset-free force rejection, and a null-space realization that preserves
higher-priority balance tasks.
We further prove that classical operational-space impedance is the
infinite-horizon, zero-input limit of the same predictive interaction law;
impedance is therefore not a separate design principle but a limiting case of
normalized predictive interaction dynamics.
Simulation on a 17-DOF biped and a Unitree G1 model validates the structural
claims: the contact-mode-indexed disturbance observer delivers offset-free force
rejection---steady-state error under a sustained 8\,N load drops from
${\sim}13$\,mm without the observer to ${\sim}0.1$\,mm with it---the
contact-consistent inertia further improves static precision over a free-space
model ($0.11$ vs $1.00$\,mm), and the constant-model predictor operates at
${\geq}1$\,kHz across genuine contact-mode switches.
The results position whole-body impedance behavior, disturbance rejection,
and contact-mode switching as consequences of a single interaction-dynamics
representation.
\end{abstract}

\begin{IEEEkeywords}
Interaction dynamics, whole-body control, model predictive control,
impedance control, floating-base robots, physical human--robot interaction,
contact-consistent dynamics, legged manipulation.
\end{IEEEkeywords}

% ----------------------------------------------------------
\section{Introduction}

Legged and floating-base robots must increasingly do more than locomote: they
must physically interact with people and their surroundings while keeping
balance. Safe interaction on such platforms is fundamentally a problem of
\emph{interaction dynamics}---the closed-loop relation between the arm
end-effector force and motion---coupled, through the contact constraints, to
the whole-body balance task.
These objectives are tightly coupled---arm motions shift the center of
mass~(CoM), changing contact force distribution, while ground reactions
propagate back through the body and appear as disturbances at the
end-effector.
Classical fixed-base impedance control~\cite{hogan1985impedance} and its
MPC extensions~\cite{cao2023passive,cao2026impedance} cannot address this
coupling because they assume the robot base is rigidly anchored.

The dominant paradigm for whole-body control of legged systems decouples
the problem into two layers: a centroidal MPC that optimizes ground
reaction forces~(GRFs) using a linearized single rigid-body dynamics~(SRBD)
model~\cite{dicarlo2018dynamic,kim2019highly}, and an inner WBC layer that
resolves these forces into joint torques via a prioritized
QP~\cite{bellicoso2016perception,koolen2016design}.
This architecture achieves remarkable locomotive agility---the MIT Cheetah
3~\cite{dicarlo2018dynamic} executes high-speed bounding and stair
climbing---but allocates 100\% of the robot's control authority to
locomotion and base-posture maintenance.
Any external arm interaction is treated as a disturbance to be suppressed,
not as a channel to be actively regulated.
A biped reaching to assist a human standing beside it, or a quadruped
manipulating a valve while maintaining stance, cannot be handled by these
frameworks with the compliance and zero-steady-state-error guarantees
required for safe pHRI.

Conversely, impedance MPC methods designed for fixed-base
manipulators~\cite{cao2026impedance,cao2023passive,haninger2022model,roveda2020model}
lack an unactuated base state, generalized-coordinate partitioning, or
contact-consistent mass inverses.
They assume an infinite-mass ground connection and cannot model the
propagation of foot contact forces to end-effector apparent inertia.
Deploying them directly on a floating-base platform produces steady-state
torque errors and potential instability during contact transitions.

The technical gap is therefore not merely the absence of another whole-body
controller. It is the absence of a representation in which interaction
dynamics remain structurally the same as robot posture, contact mode, and
apparent inertia change. This paper closes that gap through a simple
sequence: interaction dynamics are first normalized through contact
consistency, the normalized dynamics become a linear double integrator,
predictive regulation is applied to that linear interaction model, and
classical impedance appears as the infinite-horizon limiting case.
The main contributions are:

\begin{enumerate}
\item \textbf{Interaction-dynamics normalization.}
We formulate floating-base arm pHRI as an interaction-dynamics problem and
derive the contact-consistent residual plant obtained after higher-priority
balance and contact tasks are removed through the whole-body null space.

\item \textbf{Configuration-invariant predictor.}
We prove that the normalized end-effector interaction dynamics reduce to an
exact linear double integrator with constant discrete state matrix $A_d$
within each contact mode. Contact configuration and robot posture affect the
predictor only through the input matrix $B_d^{(m)}$, via the
contact-consistent task inertia $\Larm$.

\item \textbf{Predictive interaction control and impedance equivalence.}
We regulate the normalized dynamics with a finite-horizon QP and show that
classical operational-space impedance is recovered as the infinite-horizon,
zero-input limit. Thus impedance behavior is a special case of predictive
interaction dynamics rather than a separate controller family.

\item \textbf{Floating-base realization and theory validation.}
We realize the predictor inside a priority-consistent WBC stack using
contact-mode-indexed matrices, a disturbance observer, and covariance
inflation across contact changes. Simulations on a 17-DOF biped and a
Unitree G1 model validate the constant-predictor structure,
offset-free disturbance rejection, and stable operation across contact-mode
switches.
\end{enumerate}

The remainder of this paper is organized as follows.
Section~\ref{sec:related} surveys related work.
Section~\ref{sec:dynamics} formulates interaction dynamics on floating-base
robots.
Section~\ref{sec:osf} derives the contact-consistent normalization.
Section~\ref{sec:impedancempc} presents predictive regulation of the
normalized dynamics.
Section~\ref{sec:switching} treats contact-mode changes.
Section~\ref{sec:theorem} proves the impedance-equivalence result.
Section~\ref{sec:stability} analyzes stability.
Section~\ref{sec:torque} gives the whole-body realization and compares the
framework with prior work.
Section~\ref{sec:simulation} reports theory-validation simulations.
Section~\ref{sec:conclusion} concludes.

% ----------------------------------------------------------
\section{Related Work}\label{sec:related}

\subsection{Centroidal MPC for Locomotion}

The dominant approach to predictive whole-body control projects the robot's
dynamics onto its centroidal frame.
Di~Carlo et~al.~\cite{dicarlo2018dynamic} introduced the convex SRBD
formulation for MIT Cheetah~3, linearizing the rotation dynamics about a
nominal pitch-roll and treating each support foot as a rigid contact.
The resulting time-varying linear system admits a QP solution at 40\,Hz for
the outer MPC, while an inner Whole-Body Impulse Control~(WBIC) layer maps
GRFs to joint torques at 500\,Hz.
This bi-level architecture is foundational to the present work; however,
Kim et~al.~\cite{kim2019highly} dedicate both layers entirely to
locomotion---external arm forces are filtered out through centroidal inertia
assumptions and rigid GRF assignment, and no mechanism exists for compliant
manipulation.

Bellicoso et~al.~\cite{bellicoso2016perception} proposed a hierarchical WBC
scheme for the ANYmal quadruped that solves a cascaded sequence of QPs to
track base and end-effector tasks simultaneously.
While effective for continuous trotting and terrain adaptation, the WBC
stack uses classical PD task objectives with instantaneous feedback
($N{=}1$); there is no receding-horizon mechanism to predict and pre-load
corrective torque before a disturbance develops.
Sustained contact forces produce the steady-state error quantified in
\eqref{eq:sse} below.

Koolen et~al.~\cite{koolen2016design} implemented a momentum-based
controller for Atlas that regulates centroidal momentum via a QP
distributing contact forces.
The formulation shares the centroidal perspective but does not include a
predictive loop for arm impedance and relies on high-gain stiffness to
suppress interaction errors.

Sleiman et~al.~\cite{sleiman2021unified} proposed a unified MPC framework
for whole-body dynamic locomotion and manipulation on legged robots,
simultaneously optimizing body posture, gait timing, and arm end-effector
trajectories over a receding horizon.
This represents the state of the art for integrated loco-manipulation and
is the closest prior work to the present architecture.
The key structural difference is that the arm task in~\cite{sleiman2021unified}
is formulated as a fixed-base QP subproblem: the task-space inertia uses
$M^{-1}$ rather than the contact-consistent $\Mbar$, and no integrating
disturbance state is included to guarantee zero steady-state pHRI error
under sustained contact.

\subsection{Trajectory Optimization and Nonlinear MPC}

Winkler et~al.~\cite{winkler2018gait} developed TOWR, a phase-based
trajectory optimizer that simultaneously synthesizes gait sequences,
foothold locations, and full-body motions over multi-second horizons using a
nonlinear program solved by interior-point methods.
The full rigid-body model captures leg-inertia coupling that SRBD ignores,
but restricts operation to ${\sim}20$--50\,Hz and offline planning.
Real-time pHRI management at 1\,kHz is outside the scope of this approach.

Grandia et~al.~\cite{grandia2023perceptive} extended TOWR with a real-time
model-predictive framework using sequential convex approximations, achieving
20\,Hz replanning for rough terrain.
While this improves responsiveness compared to pure offline planning, it
still lacks the integrating disturbance estimation required for zero
steady-state error under sustained contact.

\subsection{Impedance MPC for Manipulation}

Force control for robotic manipulators---encompassing impedance, admittance,
and hybrid position/force strategies---is surveyed comprehensively
in~\cite{villani2016force}.
The foundational impedance control law of Hogan~\cite{hogan1985impedance}
was unified with position and torque control into a passivity-preserving
framework by Albu-Sch\"{a}ffer, Ott, and
Hirzinger~\cite{albu2007unified}, providing the theoretical basis on which
predictive extensions build.

The present paper is a direct structural extension of a line of prior
work on saturated and predictive control.
Anti-windup designs for output tracking under actuator saturation and
constant disturbances~\cite{cao2004antiwindup}, and the associated
domain-of-attraction analysis~\cite{cao2002antiwindup}, established that an
integrating disturbance channel achieves zero steady-state tracking error
for fixed-base saturated linear systems---the foundational insight carried
forward here to the floating-base, contact-switching setting via the Kalman
augmented state.
The min--max MPC formulation for LPV systems~\cite{cao2005minmax} introduced
parameter-varying input matrices with input constraints, the direct
precursor to the contact-mode-indexed $B_d^{(m)}$ library of the present
work.
Building on these fixed-base results, Cao, Cheng, and
Li~\cite{cao2023passive} introduced a passive MPC framework for pHRI on
fixed-base manipulators in which the outer MPC optimizes impedance
parameters $\{M_d,D_d\}$ over a receding horizon; passivity is enforced via
a virtual energy tank.
Because impedance parameters enter nonlinearly into the prediction matrices,
iterative solvers are required and update rates are limited to 10--30\,Hz.
The two-layer Impedance MPC for pHRI~\cite{cao2026impedance} resolved this
rate limitation on fixed-base manipulators: an analytical feedforward
cancels gravity, Coriolis, and task-space inertia, reducing the residual
plant to a double integrator with a \emph{constant} state-transition
matrix, while the MPC decision variables become corrective \emph{forces}
$\Fmpc$ and a Kalman-augmented disturbance state provides a formal
zero-steady-state-error guarantee---sub-0.05\,mm under a sustained 15\,N
force on a 7-DOF manipulator, versus 44.8\,mm for classical impedance
control.
That framework is the direct base of the present paper, which extends it
from fixed-base manipulators to floating-base humanoids: the constant-$A_d$
structure is preserved by substituting the contact-consistent mass inverse
$\Mbar$ into the feedforward linearization, the Kalman state additionally
absorbs leg-momentum variations and SRBD approximation errors, and the
contact-mode-indexed structure of~\cite{cao2005minmax} handles stance-phase
transitions while preserving the precomputed lifted-system sparsity and
rollout structure.

Haninger, Hegeler, and Peternel~\cite{haninger2022model} optimize force
references and impedance parameters jointly using stochastic MPC with
Gaussian Process models of task forces.
This provides complementary insights into uncertainty-aware impedance
shaping but does not address floating-base dynamics, underactuation, or
contact-consistent operational-space formulation.

\subsection{Operational Space Control and Floating-Base Inverse Dynamics}

Khatib~\cite{khatib1987unified} formulated the operational space control
framework, establishing task-space inertia, Coriolis compensation, and
dynamically-consistent pseudoinverses as the mathematical foundation for
task-level manipulation.
Sentis and Khatib~\cite{sentis2005synthesis} extended this to hierarchical
synthesis of whole-body behaviors, proving that priority-ordered null-space
projection guarantees non-interference between tasks---the SK05 law that
forms the backbone of Level~2 in the present architecture.
Righetti et~al.~\cite{righetti2011inverse} unified the floating-base
inverse dynamics perspective with external contact constraints, showing how
$\Mbar$ arises naturally from an orthogonal decomposition of the constrained
dynamics.
The present work builds on~\cite{khatib1987unified,sentis2005synthesis,righetti2011inverse}
by embedding a predictive MPC layer in the residual null space of the
floating-base contact-consistent hierarchy.

% ----------------------------------------------------------
\section{Interaction Dynamics on Floating-Base Robots}\label{sec:dynamics}

\subsection{Generalized Coordinates}

A floating-base robot (humanoid, quadruped) has $n$ actuated joints and a
6-DOF unactuated base~\cite{righetti2011inverse}. The generalized
coordinates are:
\begin{equation}\label{eq:gencoord}
  q = \begin{bmatrix} q_b \\ q_j \end{bmatrix} \in \mathbb{R}^{n+6},
  \quad q_b \in SE(3),\; q_j \in \mathbb{R}^n
\end{equation}
where $q_b=(p_b,R_b)$ is the base position and orientation and $q_j$ are
the $n$ joint angles. The velocity vector is:
\begin{equation}\label{eq:vel}
  \dot{q} = \begin{bmatrix} v_b \\ \dot{q}_j \end{bmatrix} \in \mathbb{R}^{n+6},
  \quad v_b = \begin{bmatrix} \dot{p}_b \\ \omega_b \end{bmatrix}
\end{equation}

\subsection{Equations of Motion}

The floating-base equations of motion~\cite{righetti2011inverse} are:
\begin{equation}\label{eq:eom}
  M(q)\ddot{q} + C(q,\dot{q})\dot{q} + G(q) = S^\top\tau + J_c^\top(q)\lambda
\end{equation}
where $M(q)\in\mathbb{R}^{(n+6)\times(n+6)}$ is the positive-definite
inertia matrix; $h=C(q,\dot{q})\dot{q}+G(q)\in\mathbb{R}^{n+6}$ collects
Coriolis, centrifugal, and gravity terms; $S=[0_{n\times6},\,I_n]$ is the
selection matrix (the base has no direct actuation); $\tau\in\mathbb{R}^n$
are commanded joint torques; $J_c(q)\in\mathbb{R}^{n_c\times(n+6)}$ is the
contact Jacobian; and $\lambda\in\mathbb{R}^{n_c}$ are ground reaction
forces~(GRFs).

Partitioning \eqref{eq:eom} into base and joint blocks:
\begin{equation}\label{eq:partition}
  \begin{bmatrix}M_b & M_{bj} \\ M_{bj}^\top & M_j\end{bmatrix}
  \begin{bmatrix}\ddot{q}_b \\ \ddot{q}_j\end{bmatrix}
  +\begin{bmatrix}h_b \\ h_j\end{bmatrix}
  =\begin{bmatrix}0 \\ \tau\end{bmatrix}
  +\begin{bmatrix}J_{c,b}^\top \\ J_{c,j}^\top\end{bmatrix}\lambda
\end{equation}

\subsection{Rigid Contact Constraints}

A rigid contact at point $i$ enforces zero contact-point
velocity: $J_{c,i}(q)\dot{q}=0$~\cite{righetti2011inverse}.
Differentiating yields the acceleration-level constraint:
\begin{equation}\label{eq:contact_acc}
  J_{c,i}\ddot{q} = -\dot{J}_{c,i}\dot{q}
\end{equation}

Stacking all contacts and combining with \eqref{eq:eom}:
\begin{equation}\label{eq:KKT}
  \begin{bmatrix}M & -J_c^\top \\ J_c & 0\end{bmatrix}
  \begin{bmatrix}\ddot{q} \\ \lambda\end{bmatrix}
  =\begin{bmatrix}S^\top\tau-h \\ \gamma_c\end{bmatrix},
  \quad \gamma_c\triangleq-\dot{J}_c\dot{q}
\end{equation}

Solving for the GRFs:
\begin{equation}\label{eq:GRF}
  \lambda = \Lambda_c(q)\bigl(-\dot{J}_c\dot{q}-J_cM^{-1}(S^\top\tau-h)\bigr)
\end{equation}
where $\Lambda_c=(J_cM^{-1}J_c^\top)^{-1}$ is the contact-space inertia.

\subsection{Contact-Consistent Mass Inverse}

Define the contact-space inertia and the corresponding constrained
inverse~\cite{khatib1987unified,righetti2011inverse}:
\begin{equation}\label{eq:Pc}
  \Lambda_c=(J_cM^{-1}J_c^\top)^{-1}
\end{equation}
\begin{equation}\label{eq:Mbar}
  \Mbar = M^{-1}-M^{-1}J_c^\top\Lambda_cJ_cM^{-1}
\end{equation}
$\Mbar$ replaces $M^{-1}$ in all operational-space formulas when contacts
are active. Equation~\eqref{eq:Mbar} is the Schur-complement inverse of the
constrained dynamics and is symmetric positive semidefinite on the admissible
acceleration subspace when $M\succ0$ and $J_c$ has full row rank. The
associated contact-consistent projector may be written in left/right forms
depending on the metric, but \eqref{eq:Mbar} is the definition used in the
task inertia $\Lambda_i=(J_i\Mbar J_i^\top)^{-1}$.

\subsection{Centroidal Dynamics}

The centroidal momentum $h_G=[k^\top,L^\top]^\top=A(q)\dot{q}\in\mathbb{R}^6$
aggregates the robot's linear and angular momentum about its
CoM~\cite{orin2013centroidal}. Differentiating:
\begin{equation}\label{eq:centroidal}
  \dot{h}_G = G_c(q)\lambda + \begin{bmatrix}mg \\ 0\end{bmatrix},
  \quad
  G_c = \begin{bmatrix}I_3 & I_3 & \cdots \\
                       (p_1{-}p_G)^\times & (p_2{-}p_G)^\times & \cdots\end{bmatrix}
\end{equation}

For the outer MPC, \eqref{eq:centroidal} is approximated by the
\textbf{SRBD} model~\cite{dicarlo2018dynamic}, treating the robot as a
lumped mass $m$ with constant inertia $I_G$. Linearizing about a nominal
orientation yields:
\begin{equation}\label{eq:SRBD}
  \dot{x}_c = A_c x_c + B_c(\{p_i\})u_c
\end{equation}
where $x_c\in\mathbb{R}^{12}$ is the centroidal state and $u_c$ collects
contact forces. The SRBD approximation error is
$O(m_{\text{leg}}/m_{\text{total}})^2$, acceptable for robots where leg
mass is below 20--30\% of total mass.

% ----------------------------------------------------------
\section{Contact-Consistent Interaction Dynamics}
\label{sec:osf}

\subsection{Task-Space Dynamics}

For a task variable $x_i=\phi_i(q)\in\mathbb{R}^{m_i}$ with Jacobian
$J_i=\partial\phi_i/\partial q$, substituting \eqref{eq:eom} into the
task-space acceleration yields the contact-consistent task-space
dynamics~\cite{khatib1987unified,righetti2011inverse}:
\begin{equation}\label{eq:taskdyn}
  \Lambda_i(q)\ddot{x}_i + \mu_i = F_i
\end{equation}
where $\Lambda_i=(J_i\Mbar J_i^\top)^{-1}$ is the task-space inertia;
$\mu_i=\bar{J}_i^\top h - \Lambda_i\dot{J}_i\dot{q}$ collects projected
Coriolis, centrifugal, and gravity terms; $\bar{J}_i=\Mbar J_i^\top\Lambda_i$
is the dynamically-consistent pseudoinverse; and $F_i$ is the commanded
operational-space force.

\subsection{Hierarchical Task Synthesis (SK05)}

For $k$ tasks in priority order (Task~1 = highest), the Sentis--Khatib
law~\cite{sentis2005synthesis} synthesizes the control torque as:
\begin{equation}\label{eq:SK05}
  \tau = J_1^\top F_1 + \bar{N}_1^\top J_2^\top F_2
       + \bar{N}_{12}^\top J_3^\top F_3 + \cdots
       + \bar{N}_{1\cdots k}^\top\tau_{\text{null}}
\end{equation}
where $\bar{N}_{1\cdots i}=\prod_{j=1}^{i}(I-\bar{J}_jJ_j)$ is the
accumulated contact-consistent null-space projector.
Under an exact model, dynamically consistent projectors, inactive actuator
saturation, and no change in the active contact/friction inequalities, the
key property $J_j\bar{N}_{1\cdots i}=0$ for all $j\leq i$ gives zero
acceleration at all higher-priority task coordinates.
Thus the arm predictive interaction layer is designed to be non-interfering
with contact maintenance and balance in the ideal hierarchy; in
implementation, saturation and contact active-set changes are handled by
conservative force limits and by the contact-mode update protocol.

Each task force $F_i$ in \eqref{eq:SK05} is typically a PD law:
\begin{equation}\label{eq:PD}
  F_i = \Lambda_i(\ddot{x}_{di}+K_{D,i}\dot{e}_i+K_{P,i}e_i)+\mu_i
\end{equation}
Under a persistent disturbance force $F_h$:
\begin{equation}\label{eq:sse}
  e_{\infty,i}
  = K_{P,i}^{-1}\Lambda_i^{-1}F_h
  = K_{x,i}^{-1}F_h \neq 0,\qquad
  K_{x,i}\triangleq\Lambda_iK_{P,i}.
\end{equation}
Here $K_{P,i}$ is the acceleration-level proportional gain in
\eqref{eq:PD}, while $K_{x,i}$ is the equivalent Cartesian stiffness in
N/m. The simulation baselines report $K_x$ directly; therefore the D1
theoretical offset under an 8\,N force and $K_x=800$\,N/m is
$8/800=10$\,mm.
This residual steady-state error is the fundamental limitation of the SK05
PD law and the primary motivation for regulating the arm task through
predictive interaction dynamics.

% ----------------------------------------------------------
\section{Predictive Regulation of Interaction Dynamics}\label{sec:impedancempc}

\subsection{Architecture Overview}

The control objectives are:
\begin{enumerate}
  \item maintain balance with contact forces inside friction cones;
  \item track an arm end-effector reference $p_d(t)$;
  \item reject pHRI forces with zero steady-state tracking error.
\end{enumerate}

These objectives are addressed by three nested control levels:

\textbf{Level~1 (Centroidal MPC, 40--100\,Hz):} plans CoM trajectory and
GRFs over a 500\,ms horizon via the SRBD model \eqref{eq:SRBD}.

\textbf{Level~2 (WBC Hierarchy, 500\,Hz):} resolves the Level~1 GRFs into
joint torques using the SK05 law \eqref{eq:SK05} for contact and balance
tasks, leaving the arm end-effector task slot open.

\textbf{Level~3 (predictive interaction regulation, $\geq$1\,kHz):} fills the arm slot with a
receding-horizon QP that predicts and rejects pHRI disturbances, replacing
the PD law \eqref{eq:PD} with a predictive force command.

\subsection{Contact-Consistent Residual Plant}

After Level~2 commits torques for contact maintenance and balance, the
residual arm end-effector dynamics are:
\begin{equation}\label{eq:plant}
  \Larm(q)\ddot{x}_{\text{arm}}+\mu_{\text{arm}} = F_{\text{arm}}+d_{\text{ext}}
\end{equation}
where $\Larm=(J_{\text{arm}}\Mbar J_{\text{arm}}^\top)^{-1}$ uses the
contact-consistent mass inverse \eqref{eq:Mbar},
$\mu_{\text{arm}}=\bar{J}_{\text{arm}}^\top h - \Larm\dot{J}_{\text{arm}}\dot{q}$,
and $d_{\text{ext}}$ is the pHRI wrench projected to arm task space.

\subsection{Contact-Consistent Feedforward and Horizon Freezing}

We execute an analytical operational-space feedforward:
\begin{subequations}\label{eq:feedforward}
\begin{align}
  F_{\text{arm}} &= \Larm(q)\ddot{p}_d + \mu_{\text{arm}} - \Fmpc \label{eq:ff_a}\\
  \tau_{\text{ff,arm}} &= S\,\Naug^\top J_{\text{arm}}^\top F_{\text{arm}} \label{eq:ff_b}
\end{align}
\end{subequations}

Level~2 updates $\Naug(q)$, $\Larm(q)$, and $\mu_{\text{arm}}$ at 500\,Hz.
Level~3 runs at $\geq$1\,kHz; on interleaved cycles the projector and
feedforward terms are held at their most recent Level~2 values.
The configuration changes by at most $\|\dot{q}\|\Delta t_2\approx0.002$\,rad
per Level~2 interval, so the frozen-matrix error is first-order small.

Substituting the feedforward \eqref{eq:ff_a} into the residual plant
\eqref{eq:plant} cancels the $\Larm\ddot p_d+\mu_{\text{arm}}$ terms, and with
the error convention $e_{\text{arm}}=x_{\text{arm}}-p_d$ (actual minus
desired) and $d(t)\triangleq\Larm^{-1}d_{\text{ext}}$ the residual tracking
error dynamics reduce to:
\begin{equation}\label{eq:errordyn}
  \ddot{e}_{\text{arm}} = -\Larm^{-1}(q)\,\Fmpc + d(t).
\end{equation}
The corrective force $\Fmpc$ thus enters with a restoring sign, so a positive
$\Fmpc$ drives $e_{\text{arm}}\to0$.

\begin{proposition}[Constant $A_d$ and Local LTI Horizon Freezing]
\label{prop:constant_Ad}
Within a fixed contact mode, by evaluating the configuration-dependent task
inertia at the current sampling instant $k$, the discrete-time state
transition matrix for $x_{e,k}=[e_{\text{arm}}^\top,\dot{e}_{\text{arm}}^\top]^\top$
reduces to a constant linear system:
\begin{equation}\label{eq:Ad}
  A_d = \begin{bmatrix}I_3 & \Delta t\,I_3 \\ 0 & I_3\end{bmatrix},
  \quad
  B_d(q_k) = \begin{bmatrix}0 \\ -\Larm^{-1}(q_k)\,\Delta t\end{bmatrix}
\end{equation}
\end{proposition}

\begin{proof}
The continuous residual error dynamics \eqref{eq:errordyn} have state matrix
$A_c=\left[\begin{smallmatrix}0&I_3\\0&0\end{smallmatrix}\right]$, which is
nilpotent ($A_c^2=0$). The matrix exponential therefore terminates at first
order, $A_d=e^{A_c\Delta t}=I+A_c\Delta t$, which is \emph{exact} and
independent of configuration---the double-integrator structure in
\eqref{eq:Ad}. Evaluating $\Larm(q_k)$ at the current sample freezes
$B_d(q_k)=[0;\,-\Larm^{-1}(q_k)\Delta t]$ over the horizon, so $(A_d,B_d(q_k))$ is
LTI within the receding window; the configuration drift
$\|\dot q\|\Delta t_2$ between Level-2 updates enters only $B_d$ and is
first-order small.
\end{proof}

The constant $A_d$ property permits the lifted rollout pattern, powers of
$A_d$, and sparsity structure of the QP to be computed offline. The numerical
input matrix $B_d(q_k)$ still depends on the frozen task inertia $\Larm(q_k)$,
so $\Gamma$, $H$, and the linear term are updated online after horizon
freezing. Thus the online computation is a small dense update on a fixed
structure, not a full nonlinear MPC rebuild.

\subsection{Receding-Horizon QP}

The input matrix indexed to contact mode $m$ is:
\begin{equation}\label{eq:Bd}
  B_{d,k}^{(m)} = \begin{bmatrix}0 \\ -(\Lambda_{\text{arm}}^{(m)}(q_k))^{-1}\Delta t\end{bmatrix}
\end{equation}

The constant $A_d$ property permits the lifted rollout pattern, powers of
$A_d$, and sparsity structure of the QP to be computed offline. The
numerical input matrix $B_{d,k}^{(m)}$ still depends on the frozen task
inertia $\Larm(q_k)$, so $\Gamma$, $H$, and the linear term are updated online
after horizon freezing.
The $N$-step prediction matrix $\Gamma_k^{(m)}$ is constructed from
$(A_d,B_{d,k}^{(m)})$ by standard lifted-system expansion.
The receding-horizon QP is:
\begin{equation}\label{eq:QP}
  \min_{U}\;\tfrac{1}{2}U^\top H_k^{(m)}U + h_k^{(m)\top}U
  \quad\text{s.t.}\quad \|\Fmpck{k}\|_\infty\leq F_{\max}
\end{equation}
with $H_k^{(m)}=\Gamma_k^{(m)\top}\bar{Q}\Gamma_k^{(m)}+\bar{R}$ and
$h_k^{(m)}=\Gamma_k^{(m)\top}\bar{Q}x_{\text{free},k}^{(m)}$.
The contact-mode index $m$ plays the role of the scheduling variable in the
LPV-MPC framework of~\cite{cao2005minmax}, while the configuration
dependence enters through the frozen $\Larm(q_k)$ inside $B_{d,k}^{(m)}$.
The box constraint $\|\Fmpc\|_\infty\leq F_{\max}$ is chosen so that arm
joint torques $\tau_{\text{arm}}=J_{\text{arm}}^\top\Fmpc$ remain below
hardware limits at all configurations within the operating
workspace~\cite{cao2002antiwindup}.
Mapping individual joint torque limits exactly into the Cartesian QP would
require a configuration-dependent constraint matrix $J_{\text{arm}}^\top$;
we therefore use the conservative Cartesian force bound to keep a fixed QP
structure even though the numerical Hessian is updated from $\Larm(q_k)$.
The QP~\eqref{eq:QP} is strictly convex and solved by OSQP~\cite{stellato2020osqp}
in $<0.1$\,ms for $N=20$.

\subsection{Kalman Disturbance Augmentation}

The disturbance $d(t)$ in \eqref{eq:errordyn} captures: (i)~external pHRI
forces; (ii)~unmodeled contact reactions from the feet; and (iii)~SRBD
approximation error.
Augmenting with an integrating disturbance state $\dhat\in\mathbb{R}^3$:
\begin{equation}\label{eq:augstate}
  \begin{bmatrix}x_{e,k+1}\\\dhat_{k+1}\end{bmatrix}
  =\underbrace{\begin{bmatrix}A_d & B_d^{(m)} \\ 0 & I\end{bmatrix}}_{A_{\text{aug}}}
  \begin{bmatrix}x_{e,k}\\\dhat_{k}\end{bmatrix}
  +\begin{bmatrix}B_d^{(m)}\\0\end{bmatrix}\Fmpck{k}+w_k
\end{equation}
with $w\sim\mathcal{N}(0,Q_w)$ and $v\sim\mathcal{N}(0,R_v)$.
Unlike the strict random-walk model of~\cite{cao2026impedance}, in which
process noise drives only the disturbance state, $Q_w$ here retains a
small block on $x_e$---two orders of magnitude below the $\dhat$
block (Section~\ref{sec:simulation})---to regularize the filter against
the frozen-matrix and SRBD residuals of the floating-base setting; the
offset-free property is unaffected, as it follows from the integrating
structure of $A_{\text{aug}}$, not from the noise partition.
The Kalman gain $K_f$ is computed offline from the steady-state discrete
algebraic Riccati equation.
The integrating structure of $A_{\text{aug}}$ guarantees $\dhat_{k}\to d$
for any bounded constant disturbance~\cite{cao2004antiwindup}.

The free-response prediction fed into the QP is:
\begin{equation}\label{eq:freeresponse}
  x_{\text{free}}^{(m)} = A_d^N x_e + \sum_{j=0}^{N-1}A_d^j B_d^{(m)}\dhat
\end{equation}
This causes the optimizer to pre-load corrective force before the
disturbance fully manifests at the end-effector.

% ----------------------------------------------------------
\section{Contact-Mode Changes in Interaction Dynamics}\label{sec:switching}

\subsection{Contact-Mode Switch Protocol}

When a foot lifts off or touches down, $J_c$ changes discontinuously;
consequently $\Mbar$, $\Larm$, and $B_d$ all jump.
At contact switch time $t_s$, the following protocol is executed in order:
\begin{enumerate}
  \item Recompute $\Mbar$ with the new $J_c$.
  \item Recompute $\Larm$; select new mode index $m_{\text{new}}$; update $B_{d,k}^{(m_{\text{new}})}$ from the frozen configuration.
  \item \textbf{Covariance inflation:} $P_{\text{aug}}\leftarrow\alpha P_{\text{aug}}$,
        $\alpha\in[3,5]$, to allow rapid re-estimation of the disturbance.
  \item \textbf{Hold $\dhat$:} do not reset to zero---balance-related
        disturbances persist across contact transitions.
\end{enumerate}
The Kalman filter re-converges within approximately $5\tau_{\text{Kalman}}$
samples, after which the disturbance estimate is again accurate.

\subsection{Contact-Mode-Indexed Matrix Library}

For a robot cycling through $K$ repeating contact modes, the contact Jacobian
row pattern, lifted-system sparsity, and powers of $A_d$ are cached for each
mode. At each QP call, the controller freezes the current configuration,
recomputes $\Larm(q_k)$ and $B_{d,k}^{(m)}$, and updates $\Gamma_k^{(m)}$ and
$H_k^{(m)}$ on the cached structure. Since $A_d$ is constant
(Proposition~\ref{prop:constant_Ad}), $A_d^N$ is precomputed once globally.
For non-periodic or highly variable contact sequences, the same lifted-system
expansion is reconstructed online; this remains tractable because only the
input blocks depend on the current contact/configuration state while the
double-integrator state transition is fixed.

% ----------------------------------------------------------
\section{Impedance as the Infinite-Horizon Limit}\label{sec:theorem}

\begin{theorem}[Impedance Equivalence of Predictive Interaction Dynamics]\label{thm:equivalence}
Under rigid contacts, fixed contact mode, and no disturbance, the
unconstrained infinite-horizon predictive interaction controller (Level~3, $N\to\infty$)
renders a classical task-space impedance law:
\begin{equation}\label{eq:impedance}
  \Larm(q)\ddot{e}+\Larm K_v\dot{e}+\Larm K_e e = F_h
\end{equation}
with effective configuration-adaptive mass $M_{d,\emph{eff}}=\Larm(q)$
depending on both arm posture and contact configuration.
\end{theorem}

\begin{proof}
In the infinite-horizon unconstrained limit with $\dhat=0$, the QP reduces to
the static discrete-LQR (Riccati) feedback $\Fmpc=K_\infty x_e$. Partition
$K_\infty=[K_p^\star\;K_v^\star]$ and define the acceleration-level gains by
factoring through the task inertia,
$K_e=\Larm^{-1}K_p^\star$ and $K_v=\Larm^{-1}K_v^\star$. Equivalently,
$\Fmpc=\Larm(q)\,(K_e\,e+K_v\,\dot e)$. Substituting into
\eqref{eq:errordyn} with
$d=\Larm^{-1}F_h$ gives
$\ddot e=-(K_e e+K_v\dot e)+\Larm^{-1}F_h$, and premultiplying by $\Larm$
yields \eqref{eq:impedance}. Hence the effective mass, damping, and stiffness
are $M_{d,\text{eff}}=\Larm(q)$, $D_{\text{eff}}=\Larm K_v$, and
$K_{\text{eff}}=\Larm K_e$.
The gains $(K_e,K_v)$ are the Riccati image of the QP weights $(Q,R)$, not
the weights themselves; no direct equality between cost weights and physical
stiffness/damping is assumed.
The contact-consistent $\Mbar$ in $\Larm$ makes the effective mass adapt to
both the arm joint configuration and the active contact footprint.
\end{proof}

Theorem~\ref{thm:equivalence} establishes that predictive interaction
regulation generalizes classical impedance control: the finite-horizon
constrained QP adds predictive disturbance rejection, constraint enforcement,
and contact-mode adaptation while reducing to a classical linear impedance
law in the infinite-horizon unconstrained limit.
The QP weights $(Q,R)$ tune the Riccati gains $(K_e,K_v)$ indirectly; the
physical stiffness and damping are the resulting
$K_{\text{eff}}=\Larm K_e$ and $D_{\text{eff}}=\Larm K_v$, not the weights
themselves.

Notably, all three closed-loop impedance parameters---$M_{d,\text{eff}}$,
$D_{\text{eff}}$, and $K_{\text{eff}}$---adapt automatically as the arm
configuration and contact state change, while the QP weights $(Q,R)$ remain
fixed design parameters.
This adaptation is a structural by-product of $\Larm(q)$ and incurs no
additional solver cost, preserving the ${\geq}1$\,kHz control rate.

% ----------------------------------------------------------
\section{Stability of Normalized Interaction Dynamics}\label{sec:stability}

\subsection{Zero Steady-State Error Under Fixed Contact Mode}

\begin{theorem}[Zero Steady-State Error]\label{thm:zss}
Suppose the disturbance $d(t)$ in \eqref{eq:errordyn} satisfies
$\|d\|_\infty\leq\bar{d}<\infty$ and is asymptotically constant.
Under fixed contact mode, the Kalman-augmented closed-loop system with the
Level~3 QP feedback is asymptotically stable, and
$\lim_{k\to\infty}\|e_{\emph{arm},k}\|=0$.
\end{theorem}

\begin{proof}
The claim combines regulator stability with offset-free disturbance rejection;
these rest on two \emph{distinct} structural properties, which the augmented
model separates.

\emph{Regulator (control) side.} The \emph{un-augmented} pair
$(A_d,B_d^{(m)})$ is controllable: $B_d^{(m)}=[0;\,-\Larm^{-1}\Delta t]$ has
full column rank whenever $\Larm^{-1}$ is nonsingular, and then
$[\,B_d^{(m)}\;\; A_dB_d^{(m)}\,]$ has rank $6$. Two singularities can violate
this---(i)~\emph{kinematic arm singularities} (e.g., a fully-extended elbow)
where $J_{\text{arm}}$ drops rank, and (ii)~\emph{contact-state singularities}
where the platform loses all ground contact---so the result is stated on a
compact, singularity-free subset $\mathcal{W}$. On $\mathcal{W}$ the
infinite-horizon LQR gain $K_\infty$ places the error-loop eigenvalues strictly
inside the unit disk.

\emph{Offset-free (observer) side.} The integrating block of
$A_{\text{aug}}$ has eigenvalues at $1$ and is \emph{not} controllable from
$\Fmpc$---the input $[B_d^{(m)};0]$ is zero on the $\dhat$ rows. This is by
design: a constant disturbance is rejected through the internal-model plus
observer, not by driving the $\dhat$ modes. The property actually required is
\emph{detectability} of $(A_{\text{aug}},C)$ with $C=[I_6\;0]$; since $\dhat$
enters the measured state $x_e$ through $B_d^{(m)}$ of full column rank, the
augmented pair is detectable, and the steady-state Kalman filter gives
$\dhat_{k}\to d$ for any bounded constant $d$. The free response
\eqref{eq:freeresponse} then pre-loads $-\dhat$, and the stable error loop
converges to $e_\infty=0$. The box constraint $\|\Fmpc\|_\infty\le F_{\max}$
defines a polyhedral invariant set $\mathcal{S}$~\cite{cao2002antiwindup} inside
which this convergence holds.
\end{proof}

\subsection{Transient Bound Across Contact Transitions}

Across a contact switch, $B_d$ jumps by
$\Delta B_d=B_d^{(m_{\text{new}})}-B_d^{(m_{\text{old}})}$.
The error bound is:
\begin{equation}\label{eq:transient}
  \|e(t)\| \leq \|e(t_s)\|
    + c_1\|\Delta B_d\|\cdot\|\dhat\|
    + c_2\|d_{\text{new}}\|\cdot\Delta t_{\text{conv}}
\end{equation}
where $\Delta t_{\text{conv}}\approx5\tau_{\text{Kalman}}$ is the
re-convergence time after covariance inflation, and $c_1,c_2>0$ depend on
the closed-loop eigenvalue placement.

\subsection{Null-Space Barrier Stability}

The contact-consistent null-space torque enforcing joint limits is:
\begin{equation}\label{eq:null}
  \tau_{\text{null}}
    = \bar{N}_{\text{arm}}^\top\bigl(-k_{\text{null}}(q-q_0)
      -d_{\text{null}}\dot{q}+g(q)\bigr)
\end{equation}
where $\bar{N}_{\text{arm}}=I-\bar{J}_{\text{arm}}J_{\text{arm}}$ uses the
contact-consistent pseudoinverse and $g(q)$ is the joint-limit barrier
gradient.
Under the same ideal hierarchy assumptions used above, projection through
$\bar{N}_{\text{arm}}$ gives zero wrench at the arm task coordinate. With
torque saturation or friction-cone active-set changes, this decoupling becomes
approximate and must be enforced by the low-level QP limits.

% ----------------------------------------------------------
\section{Whole-Body Realization and Framework Comparison}
\label{sec:torque}

\subsection{Whole-Body Torque Realization}

Combining all three levels:
\begin{equation}\label{eq:fulltorque}
  \tau = \tau_{\text{contact}}
       + \bar{N}_1^\top\tau_{\text{balance}}
       + \Naug^\top\bigl[\tau_{\text{ff,arm}}
         + J_{\text{arm}}^\top\Fmpc + \tau_{\text{null}}\bigr]
\end{equation}
Equation~\eqref{eq:fulltorque} is the whole-body realization of the
normalized interaction-dynamics controller.
The hierarchical null-space structure follows the SK05
law~\cite{sentis2005synthesis} extended to the contact-consistent
floating-base setting~\cite{khatib1987unified,righetti2011inverse}.
The contribution here is not the null-space hierarchy itself, but the
normalized predictive interaction force $\Fmpc$ that occupies the residual
arm channel. Under exact dynamics and inactive saturation/friction active-set
changes, the projections $\bar{N}_1^\top$ and $\Naug^\top$ isolate this
interaction layer from contact maintenance and balance, while the
Kalman-augmented QP provides the offset-free disturbance rejection that the
classical OS PD law \eqref{eq:PD} cannot achieve.

\subsection{Architectural Comparison}\label{sec:comparison}

Table~\ref{tab:comparison} summarizes the mathematical and architectural
distinctions between the proposed framework and the most relevant prior
work. The proposed input matrix is indexed by contact mode and by the frozen
configuration through $\Larm(q_k)$.

\begin{table*}[!t]
\caption{Architectural and Mathematical Comparison}
\label{tab:comparison}
\renewcommand{\arraystretch}{1.3}
\begin{tabular*}{\textwidth}{@{\extracolsep{\fill}}llllll}
\toprule
\textbf{Aspect} & \textbf{Base \cite{cao2026impedance}} & \textbf{Bellicoso et al.\ \cite{bellicoso2016perception}} & \textbf{Kim et al.\ \cite{kim2019highly}} & \textbf{Proposed} \\
\midrule
Primary objective    & Fixed-base pHRI           & Quadruped locomotion  & Dynamic locomotion    & Normalized interaction dynamics for floating-base pHRI \\
Task-space inertia   & $(JM^{-1}J^\top)^{-1}$   & N/A                   & N/A                   & $(J\Mbar J^\top)^{-1}$ \\
Prediction horizon   & $N$-step QP               & $N=1$                 & $\sim$10--30 steps    & $N$-step QP (${\geq}$1\,kHz) \\
Disturbance handling & Kalman, fixed base        & WBC weight tuning     & Centroidal inertia    & Kalman: pHRI + leg momentum \\
Input matrix         & $B_d$ fixed               & N/A                   & N/A                   & $B_{d,k}^{(m)}$, contact/configuration indexed \\
Null-space use       & Joint centering           & Posture tracking      & Locomotion            & Predictive interaction QP \\
Steady-state error   & Zero (fixed base)         & Nonzero under load    & Nonzero under load    & Zero (Theorem~\ref{thm:zss}) \\
Contact transitions  & N/A                       & Mode switching        & Gait phases           & Covariance-inflation protocol \\
\bottomrule
\end{tabular*}
\end{table*}

% ----------------------------------------------------------
\section{Theory-Validation Experiments}\label{sec:simulation}

The simulations are organized as validation of the theory rather than as a
catalog of scenarios. Scenarios~A and~C test the fixed-contact normalized
predictor and offset-free disturbance rejection. Scenario~B stress-tests the
same predictor under periodic contact shocks; with both feet planted at the
reported gains the disturbance observer is the dominant factor and the
contact-consistent and free-space predictors coincide. Scenarios~E and~F test
whether the same interaction layer remains well behaved when the support model
changes and the input matrix is switched through the contact-mode library. The
experiments therefore target the paper's central claims: constant normalized
state dynamics, contact-dependent input matrices, offset-free disturbance
rejection in the normalized coordinates, and stable operation across
contact-mode switches.

\subsection{Simulation Platform}

All experiments were conducted in MuJoCo~3.2~\cite{todorov2012mujoco} at a
2\,kHz integration rate.
Scenarios~A and~B use a biped comprising a 3-DOF right arm, two 4-DOF legs,
and a 6-DOF unactuated floating base (17~DOF total, 11~actuated), with
total mass 46\,kg.
Scenario~C uses the official \textbf{Unitree G1} MJCF model from MuJoCo
Menagerie (29~DOF, 33.3\,kg, factory CAD kinematics), augmented with a
single end-effector site at \texttt{right\_wrist\_yaw\_link}.
The G1 model uses position actuators ($K_p=500$, $\text{dampratio}=1$);
Level~3 applies the position-as-torque approximation
($\Delta q_i=\tau_i/K_p$) to inject predictive interaction forces through the
position channel.

The legs are held in double support by a joint-space PD balance controller;
the centroidal-MPC balance planner of Section~\ref{sec:impedancempc} is not
exercised in these static-/fixed-stance scenarios. The normalized interaction
layer runs at 1\,kHz with $N=20$, $\Delta t_3=1$\,ms.
Friction-cone half-angle $\mu=0.6$.
Level~3 cost weights: $Q=\mathrm{diag}(6\times10^4 I_3,\;60\,I_3)$,
$R=0.01\,I_3$.
Kalman noise: $Q_w=\mathrm{diag}(10^{-4}I_6,\;10^{-2}I_3)$,
$R_v=10^{-6}I_3$.
Covariance-inflation coefficient $\alpha=4$.

\subsection{Benchmarked Controllers}

Seven controllers are evaluated across all scenarios (Table~\ref{tab:controllers}).

\begin{table}[!t]
\caption{Benchmarked Controllers}
\label{tab:controllers}
\renewcommand{\arraystretch}{1.2}
\begin{tabular}{lp{6cm}}
\toprule
\textbf{Label} & \textbf{Description} \\
\midrule
D1 & Operational-space PD (task PD, no priority hierarchy), Cartesian stiffness $K_x=800$\,N/m, damping $D_x=40$\,Ns/m \\
D2 & Operational-space PI: adds $K_I=150$\,N/(m$\cdot$s) with anti-windup \\
D3 & Fixed-base impedance MPC~\cite{cao2026impedance} (uses $M^{-1}$ instead of $\Mbar$) \\
D4 & WBC assembler + null-space PD centering (no prediction) \\
D5 & Proposed: WBC + predictive interaction regulation, no Kalman \\
D6 & Proposed: WBC + predictive interaction regulation + Kalman, no inflation ($\alpha=1$) \\
D7 & Proposed full: WBC + predictive interaction regulation + Kalman + covariance inflation ($\alpha=4$) \\
\bottomrule
\end{tabular}
\end{table}

D1 establishes the analytical baseline: for Cartesian stiffness
$K_x=800$\,N/m and $F_h=8$\,N, \eqref{eq:sse} predicts
$e_\infty=F_h/K_x=10.0$\,mm exactly.

\subsection{Scenario A: Fixed Double-Support Step Disturbance}

The robot holds a stationary double-support stance.
At $t=0.5$\,s, a step pHRI force of 8\,N is applied at the end-effector in
the $x$-direction and held for the remaining 4.5\,s.
Results are shown in Table~\ref{tab:scenarioA} and Fig.~\ref{fig:scenarioA}.

\begin{table}[!t]
\caption{Scenario A --- Fixed Stance, 8\,N Step Disturbance}
\label{tab:scenarioA}
\renewcommand{\arraystretch}{1.2}
\begin{tabular}{@{}lcc@{}}
\toprule
\textbf{Controller} & \textbf{RMS err [mm]} & \textbf{SS err [mm]} \\
\midrule
D1 OS PD          & 9.24  & 10.17 \\
D2 OS PI          & 6.82  &  5.86 \\
D3 Fixed-base MPC   & 3.33  &  1.00 \\
D4 WBC + PD         & 10.06 & 10.87 \\
D5 Proposed, no Kalman & 11.82 & 12.89 \\
D6 Proposed, no inflation & \textbf{3.48} & \textbf{0.11} \\
D7 Proposed full    & \textbf{3.48} & \textbf{0.11} \\
\bottomrule
\end{tabular}
\end{table}

\begin{figure}[!t]
  \centering
  \includegraphics[width=\columnwidth]{scenario_a_results.png}
  \caption{Scenario A --- Fixed double-support stance, 8\,N step pHRI at
           $t=0.5$\,s.
           \textit{Top:} Cartesian end-effector error norm $\|e\|$ over
           time.
           The Kalman-augmented predictive controller removes most of the
           steady-state offset, while the no-Kalman predictive controller is
           biased by the sustained pHRI load.
           \textit{Bottom:} RMS and steady-state~(SS) error bar chart.
           D4--D6 are omitted for readability and reported in
           Table~\ref{tab:scenarioA}.}
  \label{fig:scenarioA}
\end{figure}

Key findings: (1)~The reactive baselines (D1, D4) leave the theoretical
$10$\,mm impedance offset ($8\,\text{N}/800\,\text{N/m}$). (2)~The disturbance
observer and input-centered QP are essential under sustained force: D6/D7
reduce the steady-state error to $0.11$\,mm, whereas the no-Kalman predictive
controller (D5) is biased to $12.9$\,mm by the unestimated load. (3)~Even in
static double support the contact-consistent inertia improves precision: the
contact-consistent D7 reaches $0.11$\,mm steady-state versus $1.00$\,mm for the
free-space fixed-base D3. Covariance inflation has no effect here (D6$=$D7)
because the contact mode never switches.

\subsection{Scenario B: Stance with Periodic Contact-Transition Shocks}

The robot holds double-support stance while a sustained 8\,N pHRI force
and periodic 6\,N spikes (1\,Hz, 0.1\,s duration) are applied.
Results are shown in Table~\ref{tab:scenarioB} and Fig.~\ref{fig:scenarioB}.

\begin{table}[!t]
\caption{Scenario B --- Stance + 1\,Hz Shocks, Sustained 8\,N pHRI}
\label{tab:scenarioB}
\renewcommand{\arraystretch}{1.2}
\begin{tabular}{lcc}
\toprule
\textbf{Controller} & \textbf{RMS err [mm]} & \textbf{Peak at switch [mm]} \\
\midrule
D1 OS PD          & 10.94 & 15.77 \\
D2 OS PI          &  5.91 & 10.17 \\
D3 Fixed-base MPC   & 3.40 & 4.84 \\
D4 WBC + PD         & 11.89 & 17.24 \\
D5 Proposed, no Kalman & 13.92 & 18.20 \\
D6 Proposed, no inflation & \textbf{3.41} & \textbf{4.85} \\
D7 Proposed full    & \textbf{3.41} & \textbf{4.85} \\
\bottomrule
\end{tabular}
\end{table}

Under sustained load plus periodic shocks, the disturbance observer is again
the dominant factor: the no-Kalman predictive controller (D5, $13.92$\,mm RMS)
is biased by the sustained load, while the Kalman-augmented controllers track
to ${\approx}3.4$\,mm. Because both feet stay planted throughout---the shocks
are applied as force spikes, not an actual contact-set change---the free-space
(D3, $3.40$\,mm) and contact-consistent (D7, $3.41$\,mm) predictors are
indistinguishable here; the value of contact-consistency requires the support
model itself to switch (Scenarios~E and~F). Covariance inflation is inert
(D6$=$D7) because the contact mode never changes.

\begin{figure}[!t]
  \centering
  \includegraphics[width=\columnwidth]{scenario_b_results.png}
  \caption{Scenario B --- double-support stance, sustained 8\,N pHRI
           + 6\,N periodic shocks at 1\,Hz.
           The disturbance observer removes the steady-state bias (D5 is
           offset, D6/D7 track to a few mm); with both feet planted, the
           free-space (D3) and contact-consistent (D7) predictors coincide
           (Table~\ref{tab:scenarioB}).}
  \label{fig:scenarioB}
\end{figure}

\subsection{Scenario C: Unitree G1 Real Model, Fixed Stance}

Scenario~A is repeated on the official Unitree G1 MJCF (29~DOF, 33.3\,kg).
Results are shown in Table~\ref{tab:scenarioC} and Fig.~\ref{fig:scenarioC}.

\begin{table}[!t]
\caption{Scenario C --- Unitree G1 (33.3\,kg), Fixed Stance, 8\,N Step}
\label{tab:scenarioC}
\renewcommand{\arraystretch}{1.2}
\begin{tabular}{lcc}
\toprule
\textbf{Controller} & \textbf{RMS err [mm]} & \textbf{SS err [mm]} \\
\midrule
D1 OS PD          & 9.00  & 9.57  \\
D2 OS PI          & 6.41  & 5.32  \\
D3 Fixed-base MPC   & 2.37  & 2.84  \\
D4 WBC + PD         & 9.00  & 9.57  \\
D5 Proposed, no Kalman & 7.54 & 8.19 \\
D6 Proposed, no inflation & \textbf{2.37} & \textbf{2.84} \\
D7 Proposed full    & \textbf{2.37} & \textbf{2.84} \\
\bottomrule
\end{tabular}
\end{table}

\begin{figure}[!t]
  \centering
  \includegraphics[width=\columnwidth]{scenario_c_g1_results.png}
  \caption{Scenario C --- Unitree G1 official model (33.3\,kg, 29~DOF),
           fixed stance, 8\,N step pHRI.
           In static G1 stance the contact-consistent (D7) and fixed-base
           (D3) predictive controllers are indistinguishable; the
           Kalman-augmented controllers improve about $3\times$ over D1.
           D4--D6 are omitted for readability and reported in
           Table~\ref{tab:scenarioC}.}
  \label{fig:scenarioC}
\end{figure}

As in Scenario~A, the static G1 stance does not distinguish contact-consistent
from fixed-base predictive control: D3 ($2.84$\,mm SS) and D7 ($2.84$\,mm) are
indistinguishable, and the Kalman-augmented predictive controllers improve
about $3\times$ over the reactive D1 ($9.57$\,mm); the no-Kalman predictive
controller (D5, $8.19$\,mm) barely improves on D1. The residual millimetre-level
error (vs.\ sub-millimetre on the biped) arises from the 5\,Hz position-actuator bandwidth
attenuating 1\,kHz MPC corrections; direct joint-torque mode on G1/R1~EDU
hardware is expected to recover near-full performance.

\subsection{Scenario E: Bracing-Hand Support Transition}

To exercise a \emph{genuine} contact-mode transition without single-support
balance, the biped is given a left arm that periodically braces against and
releases a fixed rail while both feet remain planted.
The active contact set alternates
$\{\text{L foot},\text{R foot}\}\leftrightarrow\{\text{L foot},\text{R foot},\text{L hand}\}$
($J_c$: $6\leftrightarrow9$ rows), so $\Mbar$, the right-arm task inertia
$\Larm$, and the input matrix $B_d$ genuinely switch at each brace and release.
A sustained 8\,N pHRI force acts on the right (task) arm throughout.

\begin{table}[!t]
\caption{Scenario E --- Bracing-Hand Support Transition, Sustained 8\,N pHRI}
\label{tab:scenarioE}
\renewcommand{\arraystretch}{1.2}
\footnotesize
\begin{tabular}{@{}lcc@{}}
\toprule
\textbf{Controller} & \textbf{RMS err [mm]} & \textbf{Peak at switch [mm]} \\
\midrule
D5 no Kalman              & 13.28 & 12.98 \\
D6 Kalman, no infl. ($\alpha=1$) & 4.67 & 6.69 \\
D7 Kalman + infl. ($\alpha=4$)   & \textbf{4.14} & \textbf{6.29} \\
\bottomrule
\end{tabular}
\end{table}

This scenario validates the contact-mode-indexed disturbance estimator across a
genuine contact-set switch. The no-Kalman controller is biased by the sustained
pHRI load ($13.28$\,mm RMS), while the Kalman-augmented controllers track the
switch to ${\sim}4$\,mm. Here covariance inflation gives a small further
improvement (D7 $4.14$\,mm / $6.29$\,mm peak vs D6 $4.67$ / $6.69$), consistent
with its role of re-estimating $\dhat$ under the changed $B_d^{(m)}$ after a
real contact-set change. Scenario~F extends the contact-mode
test to a scheduled \emph{single$\leftrightarrow$double} support-mode
transition for the interaction layer.

% ----------------------------------------------------------
\subsection{\texorpdfstring{Scenario F: Quasi-Static Single$\leftrightarrow$Double Support Transition}{Scenario F: Quasi-Static Single-Double Support Transition}}
\label{sec:scenarioF}
Scenarios~B and~E hold both feet planted; this scenario exercises a scheduled
change of the \emph{foot-support model} used by the interaction layer, in
which the biped shifts its weight over the right foot, lifts the left foot,
switches the contact-consistent predictor from double to right-foot support,
holds, and places the foot back:
\begin{equation*}
\{\text{L foot},\text{R foot}\}\;\longleftrightarrow\;\{\text{R foot}\},
\end{equation*}
so the contact Jacobian used by the interaction-layer model ($J_c$:
$6\leftrightarrow3$ rows), $\Mbar$, the arm task inertia $\Larm$, and the input
matrix $B_d$ all switch at the lift and the place.
The interaction layer runs above the balance controller with a sustained
$8$\,N pHRI force on the arm; the arm regulates a target fixed relative to the
torso, isolating its disturbance-rejection task from the balancing motion.

\emph{Balance stand-in and its limitations.} Realizing a standing single-support
phase requires \emph{ankle-roll} (lateral centre-of-pressure) authority, which
the biped of Scenarios~A/B/E lacks; hip-roll balancing alone produces a
foot-tipping moment that cannot be countered and the robot falls. We therefore
use a modified model (\texttt{biped\_qstatic}: a wider $18$\,cm foot and added
ankle-roll actuators) and a hand-tuned quasi-static controller---a stiff
hip-roll CoM regulator for the coarse weight transfer plus a
centre-of-pressure--limited ankle-roll torque for the fine support-mode
stabilization. This is a deliberate minimal stand-in for the Level-1 balance
layer, \emph{not} the centroidal MPC of Section~\ref{sec:impedancempc}.
MuJoCo contact auditing shows that the lifted foot can retain intermittent
toe/edge contact with the floor, so Scenario~F should be read as a
support-mode/contact-model switch test for the interaction layer, not as a full
dynamic walking or physically clean single-support result.

\begin{table}[!t]
\caption{Scenario F --- Quasi-Static Single$\leftrightarrow$Double Support
         Transition, Sustained 8\,N pHRI (torso-relative arm error)}
\label{tab:scenarioF}
\renewcommand{\arraystretch}{1.2}
\centering
\begin{tabular}{@{}lcc@{}}
\toprule
\textbf{Controller} & \textbf{RMS err [mm]} & \textbf{Peak at switch [mm]} \\
\midrule
D5 no Kalman               & 30.08 & 37.57 \\
D6 Kalman, no infl. ($\alpha=1$) & 15.95 & 24.04 \\
D7 Kalman + infl. ($\alpha=4$)   & \textbf{15.81} & \textbf{23.96} \\
\bottomrule
\end{tabular}
\end{table}

The biped stays upright across the whole cycle (minimum torso
height $0.86$\,m; scheduled single-support-mode interval ${\approx}2.1$\,s), and
the interaction layer keeps the torso-relative arm error to ${\approx}1.6$\,cm.
On this model the arm task inertia barely changes across the foot-support
switch---$\Larm$ diagonal $[1.10,1.04,2.42]$\,kg in double support versus
$[1.09,1.04,2.42]$\,kg in right-foot support ($<1\%$)---because the right arm
couples only weakly to the left-leg contact. The scenario therefore exercises
the contact-mode indexing and $B_d^{(m)}$ switching across a genuine \emph{leg}
support-mode change, but does not, on this platform, produce a large inertia
change to exploit.

The Kalman augmentation roughly halves the error, visible both in RMS
(D5 $30.08\to$ D7 $15.81$\,mm) and in the peak at the switch
(D5 $37.57\to$ D7 $23.96$\,mm). Covariance inflation is essentially neutral
(D6 $15.95$ vs D7 $15.81$\,mm): the error at a support transition is
dominated by the whole-body balancing motion---the torso adjusting under the
newly single foot---rather than by staleness of the disturbance estimate, so
faster Kalman re-convergence has limited room to help. The honest reading
across Scenarios~E and~F is that the contact-mode-indexed Kalman model is the
main source of improvement, while covariance inflation is a tunable,
scenario-dependent robustness mechanism---beneficial at the planted-foot
bracing switch (Scenario~E) but neutral here, where balancing motion dominates.

\begin{figure}[!t]
  \centering
  \includegraphics[width=\columnwidth]{scenario_f_results.png}
  \caption{Scenario F --- torso-relative arm error across a quasi-static
           scheduled single$\leftrightarrow$double support-mode transition
           (shaded: right-foot-support model). All controllers stay stable
           across the contact-model switch; the Kalman gives a modest offset
           improvement and covariance inflation is neutral
           (Table~\ref{tab:scenarioF}).}
  \label{fig:scenarioF}
\end{figure}

% ----------------------------------------------------------
\section{Conclusion}\label{sec:conclusion}

This paper proposed contact-consistent interaction dynamics normalization for
floating-base pHRI.
The central result is that, after balance and contact tasks are removed
through the whole-body hierarchy, the residual end-effector interaction
channel reduces to a linear double integrator with a constant discrete state
matrix within each contact mode.
Whole-body control, null-space projection, disturbance observation, and
contact-mode scheduling are the realization mechanisms; the underlying object
being regulated is the normalized interaction dynamics.

This viewpoint explains the empirical results.
The contact-mode-indexed disturbance observer is the primary driver of accuracy---it
delivers offset-free rejection of the sustained interaction force (steady-state
error ${\sim}0.1$\,mm vs ${\sim}13$\,mm without it)---while the
contact-consistent inertia sharpens static precision and keeps the predictor
well-conditioned across support-mode changes, because $\Larm$ captures the
contact-coupled apparent inertia
seen at the hand.
The Impedance Equivalence Theorem further shows that classical
operational-space impedance is recovered as the infinite-horizon limit of the
same predictive interaction law, so impedance becomes a design interpretation
rather than the starting point of the method.
The simulations also demonstrate operation across scheduled
support-mode/contact-model changes, although a fully dynamic single-support
walking transition with a complete centroidal-MPC balance layer remains future
work.
More broadly, the normalization perspective suggests a common foundation for
pHRI, loco-manipulation, dexterous manipulation, surgical robotics, and other
contact-rich systems: first represent the task as interaction dynamics, then
normalize those dynamics into a configuration-invariant predictive control
problem.

The proposed framework occupies a structural niche not addressed by prior
locomotion-centric frameworks~\cite{dicarlo2018dynamic,kim2019highly,bellicoso2016perception,koolen2016design}:
it deliberately halts the WBC stack after balance constraints are satisfied
and injects normalized predictive interaction regulation into the residual
null space.

Future work includes hardware validation on a Unitree R1 or G1~EDU platform
(the low-level SDK torque interface is directly compatible with the proposed
architecture), extension to variable contact modes during dynamic walking,
and integration of force-torque sensor feedback for improved Kalman
convergence at contact transitions.

% ----------------------------------------------------------
\bibliographystyle{IEEEtran}
\bibliography{references}
